Identifying compounds that act on disease-relevant proteins is central to drug discovery and repurposing, but experimentally testing every compound-target pair is impractical. Models that evaluate pairs independently do not exploit known relationships among compounds and proteins. Graph-based models capture these relationships, but extending their predictions to candidate compounds with few or no graph connections remains challenging.
To address these gaps, here we present \textbf{GIFT} (Graph-Informed Fusion for Target interaction prediction), a single inductive framework for both drug-target interaction (DTI) prediction and compound prioritization that ranks compounds for a target even when those compounds are absent from the interaction graph. GIFT combines three complementary modules. First, an edge-aware graph neural network is pretrained by contrastive learning, using hard negatives that share the compound or the protein of a true interaction, while a learnable soft-edge module proposes additional edges without directly supervising the proposed-edge scores. Second, an inductive per-node gate fuses these graph embeddings with representations learned from molecular and protein-sequence descriptors, so a prediction can be made when graph evidence is missing. Third, compound and protein nodes are randomly masked while the interaction classifier is trained, forcing it to predict interaction labels for nodes with and without graph connectivity.

We demonstrate the utility of GIFT using a lung cancer-focused knowledge graph integrating ChEMBL bioactivities, protein-protein interactions, and compound similarities. We evaluated GIFT against eight machine-learning and deep-learning baselines (random forest, GCN, GraphSAGE, GAT, DeepDTA, GraphDTA, DrugBAN, and the pretrained-representation model DTIAM) across 30 runs with different random seeds and data splits, excluding validation and test DTI edges from the message-passing graph. GIFT significantly outperformed all eight baselines (two-sided paired Wilcoxon tests, Holm-corrected) in mean PR-AUC, precision, and Matthews correlation coefficient (MCC) across three increasingly demanding benchmarks: (i) predicting held-out ChEMBL~35 interactions with random negative pairs ($p<10^{-5}$); (ii) identifying which targets the same test compounds act on using compound-matched negatives ($p<10^{-5}$); and (iii) predicting interactions recorded in ChEMBL~36 for 95 clinical-stage drugs outside the training graph ($p<10^{-7}$). These results demonstrate improved target identification and generalization beyond compounds represented during training.

The prediction advantage translated into improved compound prioritization: across 10 actionable driver kinases of non-small cell lung cancer, $4.8$ of GIFT's 10 highest-ranked drugs, on average, had known activity against the target, significantly more than DeepDTA, GraphDTA, DrugBAN, and DTIAM ($p<10^{-4}$). The top-ranked predictions included established clinically relevant lung cancer drug-target pairs, including ceritinib-ALK and dacomitinib-EGFR, demonstrating generalization to drugs outside the training graph. Demonstrated here for lung cancer, GIFT offers a reusable architecture that can be trained on compound-protein interaction data from other disease settings. By linking interaction prediction with target-specific compound ranking, GIFT provides a general framework for prioritizing candidates for more targeted experimental testing.
